%% ma: L10-B12-0019
%% lop: 10
%% chuong: 5
%% bai: 12
%% dang: 10.12.4
%% loai: TL
%% muc: TH
%% dap_an: "-"
%% nguon: "chuyen-de-cac-so-dac-trung-cua-mau-so-lieu-khong-ghep-nhom-toan-10-chuong-trinh-moi.pdf (GV. Phan Nhật Linh) (P01), Bài 01 Bài tập tự luận 19"
%% kien_thuc: "Lớp 10 Bài 12 (số gần đúng, sai số tuyệt đối, độ chính xác, sai số tương đối, quy tròn số)"
%% trang_thai: chinh-thuc
%% ghi_chu: "Lời giải gốc làm tròn số đo thành 7,4; kho tính với số đo thật, kết luận không đổi. Tự duyệt theo yêu cầu giáo viên 22/10/2026."
\begin{ex}
Học sinh đo chu kì dao động của con lắc đơn bằng đồng hồ bấm giây, bằng cách đo thời gian thực hiện một dao động toàn phần. Kết quả ba lần đo như sau (đơn vị: giây):
\begin{center}
\begin{tabular}{|c|c|c|c|}\hline
Lần đo & $1$ & $2$ & $3$\\\hline
Kết quả & $7{,}391\pm0{,}02$ & $7{,}395\pm0{,}05$ & $7{,}389\pm0{,}06$\\\hline
\end{tabular}
\end{center}
Ước lượng sai số tương đối của mỗi lần đo. Lần đo nào có sai số tương đối nhỏ nhất?
\loigiai{$\delta_1\le\dfrac{0{,}02}{7{,}391}\approx0{,}27\%$; $\delta_2\le\dfrac{0{,}05}{7{,}395}\approx0{,}68\%$; $\delta_3\le\dfrac{0{,}06}{7{,}389}\approx0{,}81\%$.

Vậy lần đo thứ nhất có sai số tương đối nhỏ nhất.}
\end{ex}
